\documentclass[10pt,a4paper]{amsart}
\usepackage[margin=19mm]{geometry}
\usepackage{amsmath,amssymb,mathtools}
\usepackage[scr=rsfs,cal=euler]{mathalfa}
\usepackage{xfrac}
\usepackage[unicode,hidelinks,bookmarks=false]{hyperref}
\newcommand{\ie}{i.e.\ }
\newcommand{\cf}{cf.\ }
\newcommand{\resp}{respectively\ }
\newcommand{\Subsection}{\subsection}
\providecommand{\nobreakdash}{\nobreak\hbox{-}}
\setlength{\emergencystretch}{2em}
\multlinegap0pt
\newcommand{\cinf}{$\mathcal{C}^{\infty}$}
\usepackage{amssymb}
\usepackage[shortlabels]{enumitem}
\usepackage{xcolor}
\usepackage{tcolorbox}
\usepackage{tikz}
\newtheorem{proposition}{Proposition}
\title{Exact integration of Hamiltonian dynamics via Jacobi and Poisson \texorpdfstring{\cinf}{cinf}-structures}
\author{A. J. Pan-Collantes, C. Sard\'on, X. Zhao}
\date{Source: arXiv:2602.11070v3 (CC BY 4.0)}
\begin{document}
\maketitle

\noindent\emph{Source and licence.} Exact integration of Hamiltonian dynamics via Jacobi and Poisson \texorpdfstring{\cinf}{cinf}-structures. Version arXiv:2602.11070v3 (CC BY 4.0), distributed by arXiv under the Creative Commons Attribution 4.0 International licence, http://creativecommons.org/licenses/by/4.0/, as recorded in the arXiv licence metadata for this exact version and shown on the abstract page https://arxiv.org/abs/2602.11070v3. Full licence notice: \texttt{licenses/arxiv-2602.11070-CC-BY-4.0.txt}.\par\smallskip

\noindent\textit{Editorial scope.} Counted unit: the article's proposition JacC (source0.tex lines 1333--1344). The article's complete original proof of it is delivered with it (source0.tex lines 1345--1358, one \texttt{proof} environment of 14 lines). No mathematics is written, repaired or re-derived: the statement and the proof are the article's own lines, in the article's own notation, and no retained line is substituted or re-flowed.  No further source-local context is needed: every cross-reference of every retained line resolves inside the two delivered spans.  Nothing was omitted from any retained span, so the package carries no omission record.  No retained line contains a citation, so the wrapper carries no bibliography and no bibliography entry.  No retained line contains a figure environment, an image-inclusion command or drawing code, so no image is delivered and the package declares no asset dependency.  Editorial formatting only, in addition: an AMS \texttt{amsart} preamble that restates the article's own declarations and macros used by the retained lines, namely the environments \texttt{proposition} on separate counters, together with the standard packages those lines need.  The retained environments print in this document as Proposition 1 in order of appearance, whereas the article prints the same items with the numbering of its own full document.  The complete native source of the article as distributed in the arXiv e-print for this exact version is bundled unchanged as \texttt{sources/194460/source0.tex}.
\begin{proposition}\label{JacC}
Let $(M,\Lambda,E)$ be a Jacobi manifold and let $f_0,\dots,f_j\in C^\infty(M)$. For any smooth function $F\in C^\infty(\mathbb{R}^{j+1})$ one has
\[
X_{F(f_0,\dots,f_j)}
=\sum_{k=0}^j \frac{\partial F}{\partial u_k}(f_0,\dots,f_j)\,X_{f_k}
+\Bigl(F(f_0,\dots,f_j)-\sum_{k=0}^j f_k\,\frac{\partial F}{\partial u_k}(f_0,\dots,f_j)\Bigr)E.
\]
In particular, if $\{f_j,f_i\}=F_{ji}(f_0,\dots,f_j)$ for some $F_{ji}$, then
\[
[X_{f_j},X_{f_i}]=X_{\{f_j,f_i\}}\in \langle X_{f_0},\dots,X_{f_j},E\rangle.
\]
\end{proposition}

\begin{proof}
By definition, $X_h=\Lambda^\sharp(dh)+hE$ for all $h\in C^\infty(M)$. Using the chain rule,
\[
d\bigl(F(f_0,\dots,f_j)\bigr)=\sum_{k=0}^j \frac{\partial F}{\partial u_k}(f_0,\dots,f_j)\,df_k.
\]
Hence
\begin{align*}
X_{F(f_0,\dots,f_j)}
&=\Lambda^\sharp\!\left(\sum_{k=0}^j \frac{\partial F}{\partial u_k}(f_0,\dots,f_j)\,df_k\right)+F(f_0,\dots,f_j)E \\
&=\sum_{k=0}^j \frac{\partial F}{\partial u_k}(f_0,\dots,f_j)\,\Lambda^\sharp(df_k)+F(f_0,\dots,f_j)E \\
&=\sum_{k=0}^j \frac{\partial F}{\partial u_k}(f_0,\dots,f_j)\,(X_{f_k}-f_kE)+F(f_0,\dots,f_j)E,
\end{align*}
which is the desired identity.
\end{proof}

\end{document}
